\section*{الفصل \ref{ch:g8:fractions} --- الكسور: العمليات الأربع}

\begin{solution}{exo:g8:fractions:1}
$\dfrac34 + \dfrac25 = \dfrac{15}{20} + \dfrac{8}{20} =
\dfrac{23}{20}$.

$\dfrac56 - \dfrac38 = \dfrac{20}{24} - \dfrac{9}{24} =
\dfrac{11}{24}$.

$\dfrac{7}{10} + \dfrac{8}{15} = \dfrac{21}{30} + \dfrac{16}{30} =
\dfrac{37}{30}$.
\end{solution}

\begin{solution}{exo:g8:fractions:2}
$-\dfrac27 + \dfrac{3}{14} = -\dfrac{4}{14} + \dfrac{3}{14} =
-\dfrac{1}{14}$.

$\dfrac16 - \dfrac54 = \dfrac{2}{12} - \dfrac{15}{12} =
-\dfrac{13}{12}$.

$-\dfrac35 - \dfrac12 = -\dfrac{6}{10} - \dfrac{5}{10} =
-\dfrac{11}{10}$.
\end{solution}

\begin{solution}{exo:g8:fractions:3}
$\dfrac58 \times \dfrac{4}{15} = \dfrac{5 \times 4}{8 \times 15}
= \dfrac{1}{2 \times 3} = \dfrac16$ (بالاختزال على $5$ وعلى $4$).

$\dfrac{9}{14} \times \dfrac73 = \dfrac{9 \times 7}{14 \times 3}
= \dfrac{3}{2}$ (بالاختزال على $7$ وعلى $3$).

$\left(-\dfrac65\right) \times \dfrac{10}{9} = -\dfrac{60}{45}
= -\dfrac43$ (إشارتان متعاكستان: سالب).
\end{solution}

\begin{solution}{exo:g8:fractions:4}
المقلوبات: $\dfrac73$؛ \quad و $\dfrac15$؛ \quad و $4$؛ \quad
و $-\dfrac92$ (\omterm{def:g8:fractions:inverse}{فالمقلوب} يحفظ الإشارة).
\end{solution}

\begin{solution}{exo:g8:fractions:5}
$\dfrac49 \div \dfrac23 = \dfrac49 \times \dfrac32 = \dfrac{12}{18}
= \dfrac23$.

$\dfrac75 \div 14 = \dfrac75 \times \dfrac{1}{14} = \dfrac{7}{70}
= \dfrac{1}{10}$.

$6 \div \dfrac34 = 6 \times \dfrac43 = 8$.
\end{solution}

\begin{solution}{exo:g8:fractions:6}
$\dfrac{3/4}{9/8} = \dfrac34 \times \dfrac89 = \dfrac{24}{36} =
\dfrac23$.

$\dfrac{5/6}{10} = \dfrac56 \times \dfrac{1}{10} = \dfrac{5}{60} =
\dfrac{1}{12}$.

$\dfrac{2}{4/7} = 2 \times \dfrac74 = \dfrac72$.
\end{solution}

\begin{solution}{exo:g8:fractions:7}
\omterm{def:g2:mult:def}{الجداء} أولًا:
$\dfrac13 + \dfrac12 \times \dfrac45 = \dfrac13 + \dfrac{4}{10}
= \dfrac{10}{30} + \dfrac{12}{30} = \dfrac{22}{30} = \dfrac{11}{15}$.

والأقواس أولًا:
$\left(\dfrac13 + \dfrac12\right) \times \dfrac45
= \dfrac56 \times \dfrac45 = \dfrac{20}{30} = \dfrac23$.
\end{solution}

\begin{solution}{exo:g8:fractions:8}
الممتلئ: $\dfrac25 + \dfrac13 = \dfrac{6}{15} + \dfrac{5}{15} =
\dfrac{11}{15}$ من الخزّان. والفارغ:
$1 - \dfrac{11}{15} = \dfrac{4}{15}$.
\end{solution}

\begin{solution}{exo:g8:fractions:9}
\omterm{def:g6:fractions:def}{الكسر} الحاصل على أكثر من $14$:
$\dfrac23 \times \dfrac34 = \dfrac{6}{12} = \dfrac12$ من القسم.
وإذا كان ذلك $12$ تلميذًا، ففي القسم $24$ تلميذًا.
\end{solution}

\begin{solution}{exo:g8:fractions:10}
$\dfrac{15}{2} \div \dfrac34 = \dfrac{15}{2} \times \dfrac43
= \dfrac{60}{6} = 10$ قطع --- عدد صحيح، فلا يضيع من الحبل شيء.
\end{solution}

\begin{solution}{exo:g8:fractions:11}
$E$: القسمة أولًا:
$\dfrac56 \div \dfrac{10}{9} = \dfrac56 \times \dfrac{9}{10}
= \dfrac{45}{60} = \dfrac34$؛ ثم
$E = \dfrac23 - \dfrac34 = \dfrac{8}{12} - \dfrac{9}{12} =
-\dfrac{1}{12}$.

$F = \left(-\dfrac12\right)^2 \times \dfrac83 = \dfrac14 \times
\dfrac83 = \dfrac{8}{12} = \dfrac23$ (فالمربع موجب).
\end{solution}

\begin{solution}{exo:g8:fractions:12}
من الداخل إلى الخارج: $1 + 1 = 2$؛ ثم
$1 + \dfrac12 = \dfrac32$؛ ثم
\[
1 + \frac{1}{3/2} = 1 + \frac23 = \frac53 .
\]
\end{solution}

\begin{solution}{pb:g8:fractions:1}
\textbf{1.} $\frac12 + \frac13 = \frac36 + \frac26 = \frac56$ و
$\frac12 + \frac14 = \frac24 + \frac14 = \frac34$. إذن
$\frac56 = \frac12 + \frac13$ و $\frac34 = \frac12 + \frac14$
كتابتان مصريتان (كسران واحديان متمايزان في كل منهما).

\textbf{2.} \omterm{def:g4:fractions:def}{بالمقام} المشترك $28$:
\[
\frac14 + \frac{1}{28} = \frac{7}{28} + \frac{1}{28}
= \frac{8}{28} = \frac{2}{7} .
\]

\textbf{3.} $\frac13 + \frac16 = \frac26 + \frac16 = \frac36 =
\frac12$، و $\frac14 + \frac{1}{12} = \frac{3}{12} +
\frac{1}{12} = \frac{4}{12} = \frac13$. وعمومًا، \omterm{def:g4:fractions:def}{المقام}
المشترك للكسرين $\frac{1}{n+1}$ و $\frac{1}{n(n+1)}$ هو
$n(n+1)$:
\[
\frac{1}{n+1} + \frac{1}{n(n+1)}
= \frac{n}{n(n+1)} + \frac{1}{n(n+1)}
= \frac{n+1}{n(n+1)}
= \frac{1}{n} .
\]

\textbf{4.} اِنطلق من $1 = \frac12 + \frac12$ وشطّر النصف
الثاني بالسؤال 3:
\[
1 = \frac12 + \frac13 + \frac16 ,
\]
أي ثلاثة كسور واحدية متمايزة.

\textbf{5.} تشطير $\frac13$ في $\frac56 = \frac12 +
\frac13$ يعطي
\[
\frac56 = \frac12 + \frac14 + \frac{1}{12} .
\]
ولا تكون أيّ كتابة مصرية وحيدة أبدًا: فمتطابقة التشطير تستطيع
دائمًا أن تعوّض \omterm{def:g6:fractions:def}{الكسر} الواحدي ذا \omterm{def:g4:fractions:def}{المقام} الأكبر
بكسرين واحديين جديدين أصغر --- فمن أيّ كتابة
يُصنع أخرى.

\textbf{6.} \omterm{def:g4:fractions:def}{بالمقام} $21$: $\frac27 = \frac{6}{21} <
\frac{7}{21} = \frac13$. وبالمقام $28$: $\frac14 =
\frac{7}{28} < \frac{8}{28} = \frac27$. والكسور الواحدية تصغر كلما
كبرت مقاماتها، $1 > \frac12 > \frac13 > \frac14 >
\dots$؛ وبما أن $\frac13$ (وكل ما فوقه) كبير جدًّا و
$\frac14$ يدخل، فإن $\frac14$ أكبر \omterm{def:g6:fractions:def}{كسر} واحدي أصغر
من $\frac27$.

\textbf{7.} $\frac27 - \frac14 = \frac{8}{28} - \frac{7}{28} =
\frac{1}{28}$: فتنتج الطريقة الجشعة
$\frac27 = \frac14 + \frac{1}{28}$، وهو بالضبط مدخل الكتبة في
السؤال 2.

\textbf{8.} \omterm{def:g4:fractions:def}{بالمقام} $14$: $\frac12 = \frac{7}{14} <
\frac{10}{14} = \frac57$، بينما \omterm{def:g6:fractions:def}{الكسر} الواحدي الوحيد الأكبر،
وهو $1$، يتجاوز $\frac57$. إذن $\frac12$ هو الأكبر الذي يدخل،
ومنه
\[
\frac57 - \frac12 = \frac{10}{14} - \frac{7}{14} = \frac{3}{14} .
\]

\textbf{9.} \omterm{def:g4:fractions:def}{بالمقام} $28$: $\frac{3}{14} = \frac{6}{28} <
\frac{7}{28} = \frac14$، إذن $\frac14$ كبير جدًّا. وبالمقام
$70$: $\frac15 = \frac{14}{70} < \frac{15}{70} =
\frac{3}{14}$، إذن $\frac15$ يدخل: فهو أكبر \omterm{def:g6:fractions:def}{كسر} واحدي
أصغر من $\frac{3}{14}$. ثم
\[
\frac{3}{14} - \frac15 = \frac{15}{70} - \frac{14}{70}
= \frac{1}{70},
\qquad\text{ومنه}\qquad
\frac57 = \frac12 + \frac15 + \frac{1}{70} .
\]
والتحقّق: $\frac{35}{70} + \frac{14}{70} + \frac{1}{70} =
\frac{50}{70} = \frac57$.

\textbf{10.} الخطوة الأولى: $\frac12 = \frac{5}{10} < \frac{8}{10} =
\frac45$ و $1 > \frac45$، إذن اِطرح $\frac12$:
$\frac45 - \frac12 = \frac{8}{10} - \frac{5}{10} = \frac{3}{10}$.
والخطوة الثانية: $\frac13 = \frac{10}{30} > \frac{9}{30} =
\frac{3}{10}$ كبير جدًّا، و $\frac14 = \frac{5}{20} <
\frac{6}{20} = \frac{3}{10}$ يدخل، إذن اِطرح $\frac14$:
$\frac{3}{10} - \frac14 = \frac{6}{20} - \frac{5}{20} =
\frac{1}{20}$. والخاتمة:
\[
\frac45 = \frac12 + \frac14 + \frac{1}{20},
\qquad
\frac{10}{20} + \frac{5}{20} + \frac{1}{20} = \frac{16}{20}
= \frac45 . \checkmark
\]

\textbf{11.} بسوط البواقي المتعاقبة هي:
من أجل $\frac27$: $2$ ثم $1$؛ ومن أجل $\frac57$: $5$ ثم $3$ ثم
$1$؛ ومن أجل $\frac45$: $4$ ثم $3$ ثم $1$. وكل قائمة
\emph{متناقصة تمامًا}.

\textbf{12.} \omterm{def:g4:fractions:def}{المقام} المشترك للكسرين $\frac ab$ و $\frac1n$
هو $b \times n$:
\[
\frac{a}{b} - \frac{1}{n}
= \frac{a \times n}{b \times n} - \frac{b}{b \times n}
= \frac{a \times n - b}{b \times n} .
\]

\textbf{13.} على \omterm{def:g4:fractions:def}{المقام} المشترك $b \times (n-1)$، تقارن
المتراجحة $\frac{1}{n-1} > \frac ab$ البسطين
$b$ و $a \times (n - 1)$:
\[
b > a \times (n - 1) = a \times n - a ,
\]
وبجمع $a - b$ إلى الطرفين نجد $a > a \times n - b$. إذن
بعد كل خطوة جشعة يكون \omterm{def:g4:fractions:def}{بسط} \omterm{def:g3:division:remainder}{الباقي} (السؤال
12) عددًا صحيحًا أصغر تمامًا من الذي قبله. والأعداد
الصحيحة لا يمكن أن تتناقص تمامًا إلى الأبد: فبعد $a$ خطوة على الأكثر
يصير \omterm{def:g4:fractions:def}{بسط} \omterm{def:g3:division:remainder}{الباقي} $1$ --- أي \omterm{def:g6:fractions:def}{كسر} واحدي --- أو
$0$، فتتوقف الطريقة.

\textbf{14.} \omterm{def:g4:fractions:def}{بالمقام} المشترك $18$:
\[
\frac12 + \frac13 + \frac19
= \frac{9}{18} + \frac{6}{18} + \frac{2}{18}
= \frac{17}{18} ,
\]
وهو \emph{أصغر من} $1$: فقد نسيت الوصية $\frac{1}{18}$ من
القطيع. وتلك الفجوة هي الحيلة كلها. فمن أصل $18$ جملًا تكون
الأنصبة $\frac12$ و $\frac13$ و $\frac19$ الأعداد الصحيحة $9$
و $6$ و $2$، ومجموعها $9 + 6 + 2 = 17$: أي القطيع بالضبط،
ويعود الجمل الثامن عشر إلى بيته. ولا يستطيع أحد أن يشتكي:
فقد وعدت الوصية، من أصل $17$ جملًا، $\frac{17}{2} = 8.5$، ثم
$\frac{17}{3} \approx 5.67$، ثم $\frac{17}{9} \approx 1.89$
جملًا، ونال كل اِبن أكثر من ذلك تمامًا ($9$ و $6$ و $2$).
\omterm{def:g6:fractions:def}{والكسر} $\frac{1}{18}$ المنسيّ هو ما ملأه جمل الجار
مؤقتًا.

\textbf{15.} يقدّم السؤال 4 الوصية المُصلَحة:
\[
\frac12 + \frac13 + \frac16 = 1 .
\]
وتكون الأنصبة أعدادًا صحيحة حين يكون \omterm{def:g6:measure:volume}{حجم} القطيع \omterm{def:g6:wholes:divisible}{قابلًا للقسمة} على
$2$ وعلى $3$ وعلى $6$ --- أي \omterm{def:g6:wholes:divisible}{قابلًا للقسمة} على $6$. وأصغر
قطيع كهذا $6$ جمال: فالأنصبة $3$ و $2$ و $1$، ولا يفضل
شيء.
\end{solution}
